\QHeading{第1题　连一连。}

\par\noindent\makebox[\linewidth][c]{
\begin{tikzpicture}[x=1bp,y=-1bp,color=black,draw=black,text=black]
\begin{scope}\clip (56.0,119.33) rectangle (545.0,266.0);
\hdhead{40}{多边形的认识——信息窗1}
\hdsection{1}{146}{连一连。}
\begin{scope}[scale=.6666667,draw=black,text=black]\draw[hdoutline] (110.000,243.000)--(146.000,188.000)--(207.000,243.000)--cycle;\end{scope}
\begin{scope}[scale=.6666667,draw=black,text=black]\draw[hdoutline] (232.000,191.000)--(265.000,240.000)--(336.000,240.000)--cycle;\end{scope}
\begin{scope}[scale=.6666667,draw=black,text=black]\draw[hdoutline] (375.000,188.000)--(444.000,188.000)--(375.000,244.000)--cycle;\end{scope}
\begin{scope}[scale=.6666667,draw=black,text=black]\draw[hdoutline] (473.000,190.000)--(473.000,241.000)--(580.000,241.000)--cycle;\end{scope}
\begin{scope}[scale=.6666667,draw=black,text=black]\draw[hdoutline] (610.000,242.000)--(635.000,188.000)--(660.000,242.000)--cycle;\end{scope}
\begin{scope}[scale=.6666667,draw=black,text=black]\draw[hdoutline] (689.000,235.000)--(753.000,194.000)--(809.000,235.000)--cycle;\end{scope}
\begin{scope}[scale=.6666667,draw=black,text=black]\draw[line width=.45pt] (383.000,188.000)--(383.000,196.000)--(375.000,196.000);\end{scope}
\begin{scope}[scale=.6666667,draw=black,text=black]\draw[line width=.45pt] (473.000,233.000)--(481.000,233.000)--(481.000,241.000);\end{scope}
\q{118}{344}{180}{13}{钝角三角形}
\q{426}{344}{185}{13}{直角三角形}
\q{685}{344}{180}{13}{锐角三角形}
\hdsection{2}{402}{分别画出下面三角形底边上的高。}
\begin{scope}[scale=.6666667,draw=black,text=black]\draw[hdoutline] (164.000,446.000)--(233.000,538.000)--(352.000,532.000)--cycle;\end{scope}
\begin{scope}[scale=.6666667,draw=black,text=black]\draw[hdoutline] (451.000,443.000)--(451.000,535.000)--(543.000,534.000)--cycle;\end{scope}
\begin{scope}[scale=.6666667,draw=black,text=black]\draw[hdoutline] (642.000,516.000)--(699.000,444.000)--(755.000,517.000)--cycle;\end{scope}
\begin{scope}[scale=.6666667,draw=black,text=black]\node[inner sep=0pt,font=\normalsize,text=black] at (273.000,474.000) {底};\end{scope}
\begin{scope}[scale=.6666667,draw=black,text=black]\node[inner sep=0pt,font=\normalsize,text=black] at (514.000,474.000) {底};\end{scope}
\begin{scope}[scale=.6666667,draw=black,text=black]\node[inner sep=0pt,font=\normalsize,text=black] at (699.000,529.000) {底};\end{scope}
\hdsection{3}{574}{画一画，标一标。}
\q{99}{620}{735}{12}{分别画一个锐角三角形和一个钝角三角形，并且标出它们的顶点、角}
\q{99}{660}{735}{12}{和边。}
\hdsection{4}{842}{按要求画一画。}
\q{101}{886}{231}{12}{① 既是直角三角形，}
\q{101}{926}{231}{12}{　又是等腰三角形。}
\q{355}{886}{233}{12}{② 既是钝角三角形，}
\q{355}{926}{233}{12}{　又是等腰三角形。}
\q{588}{886}{249}{12}{③ 既是锐角三角形，}
\q{588}{926}{249}{12}{　又是等腰三角形。}
\begin{scope}[scale=.6666667,draw=black,text=black]\draw[line width=.5pt] (130,967) rectangle (317,1154);\end{scope}
\begin{scope}[scale=.6666667,draw=black,text=black]\foreach \i in {0,...,9}{\foreach \j in {0,...,9}{\fill[black] ({143.5+\i*17.5},{979.5+\j*17.5}) circle (1.15);}}\end{scope}
\begin{scope}[scale=.6666667,draw=black,text=black]\draw[line width=.5pt] (385,967) rectangle (573,1154);\end{scope}
\begin{scope}[scale=.6666667,draw=black,text=black]\foreach \i in {0,...,9}{\foreach \j in {0,...,9}{\fill[black] ({399.5+\i*17.5},{979.5+\j*17.5}) circle (1.15);}}\end{scope}
\begin{scope}[scale=.6666667,draw=black,text=black]\draw[line width=.5pt] (618,967) rectangle (812,1154);\end{scope}
\begin{scope}[scale=.6666667,draw=black,text=black]\foreach \i in {0,...,9}{\foreach \j in {0,...,9}{\fill[black] ({632.5+\i*17.5},{979.5+\j*17.5}) circle (1.15);}}\end{scope}

\begin{scope}[color=answerred,draw=answerred,text=answerred]\draw[answerred,line width=.7pt] (103.333,166.000) -- (476.667,226.000);
\draw[answerred,line width=.7pt] (190.667,164.667) -- (112.667,226.000);
\draw[answerred,line width=.7pt] (272.000,164.000) -- (318.667,226.000);
\draw[answerred,line width=.7pt] (350.667,164.667) -- (335.333,226.000);
\draw[answerred,line width=.7pt] (423.333,164.667) -- (506.000,226.000);
\draw[answerred,line width=.7pt] (500.000,161.333) -- (134.000,226.000);
\draw[answerred,line width=.7pt,dashed] (155.333,358.667) -- (170.575,325.348);
\draw[answerred,line width=.7pt] (175.425,327.567) -- (173.206,332.417);
\draw[answerred,line width=.7pt] (173.206,332.417) -- (168.356,330.198);
\ans{166.95}{336.67}{20.00}{9}{高}
\draw[answerred,line width=.7pt,dashed] (300.667,356.667) -- (331.332,325.665);
\draw[answerred,line width=.7pt] (335.123,329.415) -- (331.373,333.207);
\draw[answerred,line width=.7pt] (331.373,333.207) -- (327.581,329.457);
\ans{320.00}{335.83}{20.00}{9}{高}
\draw[answerred,line width=.7pt,dashed] (466.000,296.000) -- (465.572,344.332);
\draw[answerred,line width=.7pt] (470.905,344.380) -- (470.953,339.047);
\draw[answerred,line width=.7pt] (470.953,339.047) -- (465.619,338.999);
\ans{469.79}{314.83}{20.00}{9}{高}
\draw[answerred,line width=.8pt] (102.000,526.667) -- (214.000,526.667) -- (152.000,467.333) -- cycle;
\draw[answerred,line width=.8pt] (321.333,526.667) -- (482.667,526.667) -- (356.667,486.667) -- cycle;
\ans{90.00}{523.33}{30.00}{9}{$\vphantom{0}A$}
\ans{216.67}{523.33}{30.00}{9}{$\vphantom{0}B$}
\ans{148.00}{454.00}{30.00}{9}{$\vphantom{0}C$}
\ans{108.67}{512.67}{30.00}{9}{$\vphantom{0}\alpha$}
\ans{194.67}{512.67}{30.00}{9}{$\vphantom{0}\beta$}
\ans{148.00}{481.33}{30.00}{9}{$\vphantom{0}\gamma$}
\ans{134.00}{528.00}{30.00}{9}{$\vphantom{0}AB$}
\ans{109.33}{490.00}{30.00}{9}{$\vphantom{0}CA$}
\ans{188.00}{490.00}{30.00}{9}{$\vphantom{0}BC$}
\ans{310.00}{523.33}{30.00}{9}{$\vphantom{0}A$}
\ans{484.67}{523.33}{30.00}{9}{$\vphantom{0}B$}
\ans{351.33}{474.00}{30.00}{9}{$\vphantom{0}C$}
\ans{338.00}{514.00}{30.00}{9}{$\vphantom{0}\alpha$}
\ans{451.33}{514.00}{30.00}{9}{$\vphantom{0}\beta$}
\ans{360.67}{496.67}{30.00}{9}{$\vphantom{0}\gamma$}
\ans{390.00}{528.00}{30.00}{9}{$\vphantom{0}AB$}
\ans{324.67}{494.00}{30.00}{9}{$\vphantom{0}CA$}
\ans{415.33}{496.67}{30.00}{9}{$\vphantom{0}BC$}
\ans{65.33}{544.67}{493.33}{8.5}{顶点：$A,B,C$；边：$AB,BC,CA$；角：$\alpha=\angle A,\ \beta=\angle B,\ \gamma=\angle C$。}
\draw[answerred,line width=.8pt] (107.333,746.667) -- (107.333,676.667) -- (177.333,746.667) -- cycle;
\draw[answerred,line width=.7pt] (107.333,740.000) -- (114.000,740.000);
\draw[answerred,line width=.7pt] (114.000,740.000) -- (114.000,746.667);
\draw[answerred,line width=.8pt] (278.000,734.667) -- (371.333,734.667) -- (324.667,711.333) -- cycle;
\draw[answerred,line width=.8pt] (433.333,746.667) -- (526.667,746.667) -- (480.000,664.667) -- cycle;
\end{scope}
\end{scope}\pgfresetboundingbox\path[use as bounding box] (56.0,119.33) rectangle (545.0,266.0);
\end{tikzpicture}}\par
\par\noindent




\par

\par\noindent

\QHeading{第2题　分别画出下面三角形底边上的高。}

\par\noindent\makebox[\linewidth][c]{
\begin{tikzpicture}[x=1bp,y=-1bp,color=black,draw=black,text=black]
\begin{scope}\clip (56.0,290.0) rectangle (509.0,380.67);
\hdhead{40}{多边形的认识——信息窗1}
\hdsection{1}{146}{连一连。}
\begin{scope}[scale=.6666667,draw=black,text=black]\draw[hdoutline] (110.000,243.000)--(146.000,188.000)--(207.000,243.000)--cycle;\end{scope}
\begin{scope}[scale=.6666667,draw=black,text=black]\draw[hdoutline] (232.000,191.000)--(265.000,240.000)--(336.000,240.000)--cycle;\end{scope}
\begin{scope}[scale=.6666667,draw=black,text=black]\draw[hdoutline] (375.000,188.000)--(444.000,188.000)--(375.000,244.000)--cycle;\end{scope}
\begin{scope}[scale=.6666667,draw=black,text=black]\draw[hdoutline] (473.000,190.000)--(473.000,241.000)--(580.000,241.000)--cycle;\end{scope}
\begin{scope}[scale=.6666667,draw=black,text=black]\draw[hdoutline] (610.000,242.000)--(635.000,188.000)--(660.000,242.000)--cycle;\end{scope}
\begin{scope}[scale=.6666667,draw=black,text=black]\draw[hdoutline] (689.000,235.000)--(753.000,194.000)--(809.000,235.000)--cycle;\end{scope}
\begin{scope}[scale=.6666667,draw=black,text=black]\draw[line width=.45pt] (383.000,188.000)--(383.000,196.000)--(375.000,196.000);\end{scope}
\begin{scope}[scale=.6666667,draw=black,text=black]\draw[line width=.45pt] (473.000,233.000)--(481.000,233.000)--(481.000,241.000);\end{scope}
\q{118}{344}{180}{13}{钝角三角形}
\q{426}{344}{185}{13}{直角三角形}
\q{685}{344}{180}{13}{锐角三角形}
\hdsection{2}{402}{分别画出下面三角形底边上的高。}
\begin{scope}[scale=.6666667,draw=black,text=black]\draw[hdoutline] (164.000,446.000)--(233.000,538.000)--(352.000,532.000)--cycle;\end{scope}
\begin{scope}[scale=.6666667,draw=black,text=black]\draw[hdoutline] (451.000,443.000)--(451.000,535.000)--(543.000,534.000)--cycle;\end{scope}
\begin{scope}[scale=.6666667,draw=black,text=black]\draw[hdoutline] (642.000,516.000)--(699.000,444.000)--(755.000,517.000)--cycle;\end{scope}
\begin{scope}[scale=.6666667,draw=black,text=black]\node[inner sep=0pt,font=\normalsize,text=black] at (273.000,474.000) {底};\end{scope}
\begin{scope}[scale=.6666667,draw=black,text=black]\node[inner sep=0pt,font=\normalsize,text=black] at (514.000,474.000) {底};\end{scope}
\begin{scope}[scale=.6666667,draw=black,text=black]\node[inner sep=0pt,font=\normalsize,text=black] at (699.000,529.000) {底};\end{scope}
\hdsection{3}{574}{画一画，标一标。}
\q{99}{620}{735}{12}{分别画一个锐角三角形和一个钝角三角形，并且标出它们的顶点、角}
\q{99}{660}{735}{12}{和边。}
\hdsection{4}{842}{按要求画一画。}
\q{101}{886}{231}{12}{① 既是直角三角形，}
\q{101}{926}{231}{12}{　又是等腰三角形。}
\q{355}{886}{233}{12}{② 既是钝角三角形，}
\q{355}{926}{233}{12}{　又是等腰三角形。}
\q{588}{886}{249}{12}{③ 既是锐角三角形，}
\q{588}{926}{249}{12}{　又是等腰三角形。}
\begin{scope}[scale=.6666667,draw=black,text=black]\draw[line width=.5pt] (130,967) rectangle (317,1154);\end{scope}
\begin{scope}[scale=.6666667,draw=black,text=black]\foreach \i in {0,...,9}{\foreach \j in {0,...,9}{\fill[black] ({143.5+\i*17.5},{979.5+\j*17.5}) circle (1.15);}}\end{scope}
\begin{scope}[scale=.6666667,draw=black,text=black]\draw[line width=.5pt] (385,967) rectangle (573,1154);\end{scope}
\begin{scope}[scale=.6666667,draw=black,text=black]\foreach \i in {0,...,9}{\foreach \j in {0,...,9}{\fill[black] ({399.5+\i*17.5},{979.5+\j*17.5}) circle (1.15);}}\end{scope}
\begin{scope}[scale=.6666667,draw=black,text=black]\draw[line width=.5pt] (618,967) rectangle (812,1154);\end{scope}
\begin{scope}[scale=.6666667,draw=black,text=black]\foreach \i in {0,...,9}{\foreach \j in {0,...,9}{\fill[black] ({632.5+\i*17.5},{979.5+\j*17.5}) circle (1.15);}}\end{scope}

\begin{scope}[color=answerred,draw=answerred,text=answerred]\draw[answerred,line width=.7pt] (103.333,166.000) -- (476.667,226.000);
\draw[answerred,line width=.7pt] (190.667,164.667) -- (112.667,226.000);
\draw[answerred,line width=.7pt] (272.000,164.000) -- (318.667,226.000);
\draw[answerred,line width=.7pt] (350.667,164.667) -- (335.333,226.000);
\draw[answerred,line width=.7pt] (423.333,164.667) -- (506.000,226.000);
\draw[answerred,line width=.7pt] (500.000,161.333) -- (134.000,226.000);
\draw[answerred,line width=.7pt,dashed] (155.333,358.667) -- (170.575,325.348);
\draw[answerred,line width=.7pt] (175.425,327.567) -- (173.206,332.417);
\draw[answerred,line width=.7pt] (173.206,332.417) -- (168.356,330.198);
\ans{166.95}{336.67}{20.00}{9}{高}
\draw[answerred,line width=.7pt,dashed] (300.667,356.667) -- (331.332,325.665);
\draw[answerred,line width=.7pt] (335.123,329.415) -- (331.373,333.207);
\draw[answerred,line width=.7pt] (331.373,333.207) -- (327.581,329.457);
\ans{320.00}{335.83}{20.00}{9}{高}
\draw[answerred,line width=.7pt,dashed] (466.000,296.000) -- (465.572,344.332);
\draw[answerred,line width=.7pt] (470.905,344.380) -- (470.953,339.047);
\draw[answerred,line width=.7pt] (470.953,339.047) -- (465.619,338.999);
\ans{469.79}{314.83}{20.00}{9}{高}
\draw[answerred,line width=.8pt] (102.000,526.667) -- (214.000,526.667) -- (152.000,467.333) -- cycle;
\draw[answerred,line width=.8pt] (321.333,526.667) -- (482.667,526.667) -- (356.667,486.667) -- cycle;
\ans{90.00}{523.33}{30.00}{9}{$\vphantom{0}A$}
\ans{216.67}{523.33}{30.00}{9}{$\vphantom{0}B$}
\ans{148.00}{454.00}{30.00}{9}{$\vphantom{0}C$}
\ans{108.67}{512.67}{30.00}{9}{$\vphantom{0}\alpha$}
\ans{194.67}{512.67}{30.00}{9}{$\vphantom{0}\beta$}
\ans{148.00}{481.33}{30.00}{9}{$\vphantom{0}\gamma$}
\ans{134.00}{528.00}{30.00}{9}{$\vphantom{0}AB$}
\ans{109.33}{490.00}{30.00}{9}{$\vphantom{0}CA$}
\ans{188.00}{490.00}{30.00}{9}{$\vphantom{0}BC$}
\ans{310.00}{523.33}{30.00}{9}{$\vphantom{0}A$}
\ans{484.67}{523.33}{30.00}{9}{$\vphantom{0}B$}
\ans{351.33}{474.00}{30.00}{9}{$\vphantom{0}C$}
\ans{338.00}{514.00}{30.00}{9}{$\vphantom{0}\alpha$}
\ans{451.33}{514.00}{30.00}{9}{$\vphantom{0}\beta$}
\ans{360.67}{496.67}{30.00}{9}{$\vphantom{0}\gamma$}
\ans{390.00}{528.00}{30.00}{9}{$\vphantom{0}AB$}
\ans{324.67}{494.00}{30.00}{9}{$\vphantom{0}CA$}
\ans{415.33}{496.67}{30.00}{9}{$\vphantom{0}BC$}
\ans{65.33}{544.67}{493.33}{8.5}{顶点：$A,B,C$；边：$AB,BC,CA$；角：$\alpha=\angle A,\ \beta=\angle B,\ \gamma=\angle C$。}
\draw[answerred,line width=.8pt] (107.333,746.667) -- (107.333,676.667) -- (177.333,746.667) -- cycle;
\draw[answerred,line width=.7pt] (107.333,740.000) -- (114.000,740.000);
\draw[answerred,line width=.7pt] (114.000,740.000) -- (114.000,746.667);
\draw[answerred,line width=.8pt] (278.000,734.667) -- (371.333,734.667) -- (324.667,711.333) -- cycle;
\draw[answerred,line width=.8pt] (433.333,746.667) -- (526.667,746.667) -- (480.000,664.667) -- cycle;
\end{scope}
\end{scope}\pgfresetboundingbox\path[use as bounding box] (56.0,290.0) rectangle (509.0,380.67);
\end{tikzpicture}}\par



\par





\par\noindent

\QHeading{第2题　作高的两种工具方法}

\Answer{方法一：三角板的一条直角边贴底边或底边的延长线，移动到另一条直角边经过对面顶点，沿该边画虚线，在虚线旁标“高”，在垂足处标直角符号。方法二：用量角器在底边所在直线上确定90度方向，使所作直线经过对面顶点，再标出垂足、虚线和“高”。两种作法得到同一条高，直角标记保持实线。}

\QHeading{第3题　画一画，标一标。}

\par\noindent\makebox[\linewidth][c]{
\begin{tikzpicture}[x=1bp,y=-1bp,color=black,draw=black,text=black]
\begin{scope}\clip (56.0,435.0) rectangle (497.0,559.33);
\hdhead{40}{多边形的认识——信息窗1}
\hdsection{1}{146}{连一连。}
\begin{scope}[scale=.6666667,draw=black,text=black]\draw[hdoutline] (110.000,243.000)--(146.000,188.000)--(207.000,243.000)--cycle;\end{scope}
\begin{scope}[scale=.6666667,draw=black,text=black]\draw[hdoutline] (232.000,191.000)--(265.000,240.000)--(336.000,240.000)--cycle;\end{scope}
\begin{scope}[scale=.6666667,draw=black,text=black]\draw[hdoutline] (375.000,188.000)--(444.000,188.000)--(375.000,244.000)--cycle;\end{scope}
\begin{scope}[scale=.6666667,draw=black,text=black]\draw[hdoutline] (473.000,190.000)--(473.000,241.000)--(580.000,241.000)--cycle;\end{scope}
\begin{scope}[scale=.6666667,draw=black,text=black]\draw[hdoutline] (610.000,242.000)--(635.000,188.000)--(660.000,242.000)--cycle;\end{scope}
\begin{scope}[scale=.6666667,draw=black,text=black]\draw[hdoutline] (689.000,235.000)--(753.000,194.000)--(809.000,235.000)--cycle;\end{scope}
\begin{scope}[scale=.6666667,draw=black,text=black]\draw[line width=.45pt] (383.000,188.000)--(383.000,196.000)--(375.000,196.000);\end{scope}
\begin{scope}[scale=.6666667,draw=black,text=black]\draw[line width=.45pt] (473.000,233.000)--(481.000,233.000)--(481.000,241.000);\end{scope}
\q{118}{344}{180}{13}{钝角三角形}
\q{426}{344}{185}{13}{直角三角形}
\q{685}{344}{180}{13}{锐角三角形}
\hdsection{2}{402}{分别画出下面三角形底边上的高。}
\begin{scope}[scale=.6666667,draw=black,text=black]\draw[hdoutline] (164.000,446.000)--(233.000,538.000)--(352.000,532.000)--cycle;\end{scope}
\begin{scope}[scale=.6666667,draw=black,text=black]\draw[hdoutline] (451.000,443.000)--(451.000,535.000)--(543.000,534.000)--cycle;\end{scope}
\begin{scope}[scale=.6666667,draw=black,text=black]\draw[hdoutline] (642.000,516.000)--(699.000,444.000)--(755.000,517.000)--cycle;\end{scope}
\begin{scope}[scale=.6666667,draw=black,text=black]\node[inner sep=0pt,font=\normalsize,text=black] at (273.000,474.000) {底};\end{scope}
\begin{scope}[scale=.6666667,draw=black,text=black]\node[inner sep=0pt,font=\normalsize,text=black] at (514.000,474.000) {底};\end{scope}
\begin{scope}[scale=.6666667,draw=black,text=black]\node[inner sep=0pt,font=\normalsize,text=black] at (699.000,529.000) {底};\end{scope}
\hdsection{3}{574}{画一画，标一标。}

\q{99}{660}{735}{12}{和边。}
\hdsection{4}{842}{按要求画一画。}
\q{101}{886}{231}{12}{① 既是直角三角形，}
\q{101}{926}{231}{12}{　又是等腰三角形。}
\q{355}{886}{233}{12}{② 既是钝角三角形，}
\q{355}{926}{233}{12}{　又是等腰三角形。}
\q{588}{886}{249}{12}{③ 既是锐角三角形，}
\q{588}{926}{249}{12}{　又是等腰三角形。}
\begin{scope}[scale=.6666667,draw=black,text=black]\draw[line width=.5pt] (130,967) rectangle (317,1154);\end{scope}
\begin{scope}[scale=.6666667,draw=black,text=black]\foreach \i in {0,...,9}{\foreach \j in {0,...,9}{\fill[black] ({143.5+\i*17.5},{979.5+\j*17.5}) circle (1.15);}}\end{scope}
\begin{scope}[scale=.6666667,draw=black,text=black]\draw[line width=.5pt] (385,967) rectangle (573,1154);\end{scope}
\begin{scope}[scale=.6666667,draw=black,text=black]\foreach \i in {0,...,9}{\foreach \j in {0,...,9}{\fill[black] ({399.5+\i*17.5},{979.5+\j*17.5}) circle (1.15);}}\end{scope}
\begin{scope}[scale=.6666667,draw=black,text=black]\draw[line width=.5pt] (618,967) rectangle (812,1154);\end{scope}
\begin{scope}[scale=.6666667,draw=black,text=black]\foreach \i in {0,...,9}{\foreach \j in {0,...,9}{\fill[black] ({632.5+\i*17.5},{979.5+\j*17.5}) circle (1.15);}}\end{scope}

\begin{scope}[color=answerred,draw=answerred,text=answerred]\draw[answerred,line width=.7pt] (103.333,166.000) -- (476.667,226.000);
\draw[answerred,line width=.7pt] (190.667,164.667) -- (112.667,226.000);
\draw[answerred,line width=.7pt] (272.000,164.000) -- (318.667,226.000);
\draw[answerred,line width=.7pt] (350.667,164.667) -- (335.333,226.000);
\draw[answerred,line width=.7pt] (423.333,164.667) -- (506.000,226.000);
\draw[answerred,line width=.7pt] (500.000,161.333) -- (134.000,226.000);
\draw[answerred,line width=.7pt,dashed] (155.333,358.667) -- (170.575,325.348);
\draw[answerred,line width=.7pt] (175.425,327.567) -- (173.206,332.417);
\draw[answerred,line width=.7pt] (173.206,332.417) -- (168.356,330.198);
\ans{166.95}{336.67}{20.00}{9}{高}
\draw[answerred,line width=.7pt,dashed] (300.667,356.667) -- (331.332,325.665);
\draw[answerred,line width=.7pt] (335.123,329.415) -- (331.373,333.207);
\draw[answerred,line width=.7pt] (331.373,333.207) -- (327.581,329.457);
\ans{320.00}{335.83}{20.00}{9}{高}
\draw[answerred,line width=.7pt,dashed] (466.000,296.000) -- (465.572,344.332);
\draw[answerred,line width=.7pt] (470.905,344.380) -- (470.953,339.047);
\draw[answerred,line width=.7pt] (470.953,339.047) -- (465.619,338.999);
\ans{469.79}{314.83}{20.00}{9}{高}
\draw[answerred,line width=.8pt] (102.000,526.667) -- (214.000,526.667) -- (152.000,467.333) -- cycle;
\draw[answerred,line width=.8pt] (321.333,526.667) -- (482.667,526.667) -- (356.667,486.667) -- cycle;
\ans{90.00}{523.33}{30.00}{9}{$\vphantom{0}A$}
\ans{216.67}{523.33}{30.00}{9}{$\vphantom{0}B$}
\ans{148.00}{454.00}{30.00}{9}{$\vphantom{0}C$}
\ans{108.67}{512.67}{30.00}{9}{$\vphantom{0}\alpha$}
\ans{194.67}{512.67}{30.00}{9}{$\vphantom{0}\beta$}
\ans{148.00}{481.33}{30.00}{9}{$\vphantom{0}\gamma$}
\ans{134.00}{528.00}{30.00}{9}{$\vphantom{0}AB$}
\ans{109.33}{490.00}{30.00}{9}{$\vphantom{0}CA$}
\ans{188.00}{490.00}{30.00}{9}{$\vphantom{0}BC$}
\ans{310.00}{523.33}{30.00}{9}{$\vphantom{0}A$}
\ans{484.67}{523.33}{30.00}{9}{$\vphantom{0}B$}
\ans{351.33}{474.00}{30.00}{9}{$\vphantom{0}C$}
\ans{338.00}{514.00}{30.00}{9}{$\vphantom{0}\alpha$}
\ans{451.33}{514.00}{30.00}{9}{$\vphantom{0}\beta$}
\ans{360.67}{496.67}{30.00}{9}{$\vphantom{0}\gamma$}
\ans{390.00}{528.00}{30.00}{9}{$\vphantom{0}AB$}
\ans{324.67}{494.00}{30.00}{9}{$\vphantom{0}CA$}
\ans{415.33}{496.67}{30.00}{9}{$\vphantom{0}BC$}
\ans{65.33}{544.67}{493.33}{8.5}{顶点：$A,B,C$；边：$AB,BC,CA$；角：$\alpha=\angle A,\ \beta=\angle B,\ \gamma=\angle C$。}
\draw[answerred,line width=.8pt] (107.333,746.667) -- (107.333,676.667) -- (177.333,746.667) -- cycle;
\draw[answerred,line width=.7pt] (107.333,740.000) -- (114.000,740.000);
\draw[answerred,line width=.7pt] (114.000,740.000) -- (114.000,746.667);
\draw[answerred,line width=.8pt] (278.000,734.667) -- (371.333,734.667) -- (324.667,711.333) -- cycle;
\draw[answerred,line width=.8pt] (433.333,746.667) -- (526.667,746.667) -- (480.000,664.667) -- cycle;
\end{scope}
\end{scope}\pgfresetboundingbox\path[use as bounding box] (56.0,435.0) rectangle (497.0,559.33);
\end{tikzpicture}}\par

\PlainItem{分别画一个锐角三角形和一个钝角三角形，并且标出它们的顶点、角}
\par





\par\noindent

\QHeading{第4题　按要求画一画。}

\par\noindent\makebox[\linewidth][c]{
\begin{tikzpicture}[x=1bp,y=-1bp,color=black,draw=black,text=black]
\begin{scope}\clip (62.0,585.0) rectangle (547.0,775.0);
\hdhead{40}{多边形的认识——信息窗1}
\hdsection{1}{146}{连一连。}
\begin{scope}[scale=.6666667,draw=black,text=black]\draw[hdoutline] (110.000,243.000)--(146.000,188.000)--(207.000,243.000)--cycle;\end{scope}
\begin{scope}[scale=.6666667,draw=black,text=black]\draw[hdoutline] (232.000,191.000)--(265.000,240.000)--(336.000,240.000)--cycle;\end{scope}
\begin{scope}[scale=.6666667,draw=black,text=black]\draw[hdoutline] (375.000,188.000)--(444.000,188.000)--(375.000,244.000)--cycle;\end{scope}
\begin{scope}[scale=.6666667,draw=black,text=black]\draw[hdoutline] (473.000,190.000)--(473.000,241.000)--(580.000,241.000)--cycle;\end{scope}
\begin{scope}[scale=.6666667,draw=black,text=black]\draw[hdoutline] (610.000,242.000)--(635.000,188.000)--(660.000,242.000)--cycle;\end{scope}
\begin{scope}[scale=.6666667,draw=black,text=black]\draw[hdoutline] (689.000,235.000)--(753.000,194.000)--(809.000,235.000)--cycle;\end{scope}
\begin{scope}[scale=.6666667,draw=black,text=black]\draw[line width=.45pt] (383.000,188.000)--(383.000,196.000)--(375.000,196.000);\end{scope}
\begin{scope}[scale=.6666667,draw=black,text=black]\draw[line width=.45pt] (473.000,233.000)--(481.000,233.000)--(481.000,241.000);\end{scope}
\q{118}{344}{180}{13}{钝角三角形}
\q{426}{344}{185}{13}{直角三角形}
\q{685}{344}{180}{13}{锐角三角形}
\hdsection{2}{402}{分别画出下面三角形底边上的高。}
\begin{scope}[scale=.6666667,draw=black,text=black]\draw[hdoutline] (164.000,446.000)--(233.000,538.000)--(352.000,532.000)--cycle;\end{scope}
\begin{scope}[scale=.6666667,draw=black,text=black]\draw[hdoutline] (451.000,443.000)--(451.000,535.000)--(543.000,534.000)--cycle;\end{scope}
\begin{scope}[scale=.6666667,draw=black,text=black]\draw[hdoutline] (642.000,516.000)--(699.000,444.000)--(755.000,517.000)--cycle;\end{scope}
\begin{scope}[scale=.6666667,draw=black,text=black]\node[inner sep=0pt,font=\normalsize,text=black] at (273.000,474.000) {底};\end{scope}
\begin{scope}[scale=.6666667,draw=black,text=black]\node[inner sep=0pt,font=\normalsize,text=black] at (514.000,474.000) {底};\end{scope}
\begin{scope}[scale=.6666667,draw=black,text=black]\node[inner sep=0pt,font=\normalsize,text=black] at (699.000,529.000) {底};\end{scope}
\hdsection{3}{574}{画一画，标一标。}
\q{99}{620}{735}{12}{分别画一个锐角三角形和一个钝角三角形，并且标出它们的顶点、角}
\q{99}{660}{735}{12}{和边。}
\hdsection{4}{842}{按要求画一画。}
\q{101}{886}{231}{12}{① 既是直角三角形，}
\q{101}{926}{231}{12}{　又是等腰三角形。}
\q{355}{886}{233}{12}{② 既是钝角三角形，}
\q{355}{926}{233}{12}{　又是等腰三角形。}
\q{588}{886}{249}{12}{③ 既是锐角三角形，}
\q{588}{926}{249}{12}{　又是等腰三角形。}
\begin{scope}[scale=.6666667,draw=black,text=black]\draw[line width=.5pt] (130,967) rectangle (317,1154);\end{scope}
\begin{scope}[scale=.6666667,draw=black,text=black]\foreach \i in {0,...,9}{\foreach \j in {0,...,9}{\fill[black] ({143.5+\i*17.5},{979.5+\j*17.5}) circle (1.15);}}\end{scope}
\begin{scope}[scale=.6666667,draw=black,text=black]\draw[line width=.5pt] (385,967) rectangle (573,1154);\end{scope}
\begin{scope}[scale=.6666667,draw=black,text=black]\foreach \i in {0,...,9}{\foreach \j in {0,...,9}{\fill[black] ({399.5+\i*17.5},{979.5+\j*17.5}) circle (1.15);}}\end{scope}
\begin{scope}[scale=.6666667,draw=black,text=black]\draw[line width=.5pt] (618,967) rectangle (812,1154);\end{scope}
\begin{scope}[scale=.6666667,draw=black,text=black]\foreach \i in {0,...,9}{\foreach \j in {0,...,9}{\fill[black] ({632.5+\i*17.5},{979.5+\j*17.5}) circle (1.15);}}\end{scope}

\begin{scope}[color=answerred,draw=answerred,text=answerred]\draw[answerred,line width=.7pt] (103.333,166.000) -- (476.667,226.000);
\draw[answerred,line width=.7pt] (190.667,164.667) -- (112.667,226.000);
\draw[answerred,line width=.7pt] (272.000,164.000) -- (318.667,226.000);
\draw[answerred,line width=.7pt] (350.667,164.667) -- (335.333,226.000);
\draw[answerred,line width=.7pt] (423.333,164.667) -- (506.000,226.000);
\draw[answerred,line width=.7pt] (500.000,161.333) -- (134.000,226.000);
\draw[answerred,line width=.7pt,dashed] (155.333,358.667) -- (170.575,325.348);
\draw[answerred,line width=.7pt] (175.425,327.567) -- (173.206,332.417);
\draw[answerred,line width=.7pt] (173.206,332.417) -- (168.356,330.198);
\ans{166.95}{336.67}{20.00}{9}{高}
\draw[answerred,line width=.7pt,dashed] (300.667,356.667) -- (331.332,325.665);
\draw[answerred,line width=.7pt] (335.123,329.415) -- (331.373,333.207);
\draw[answerred,line width=.7pt] (331.373,333.207) -- (327.581,329.457);
\ans{320.00}{335.83}{20.00}{9}{高}
\draw[answerred,line width=.7pt,dashed] (466.000,296.000) -- (465.572,344.332);
\draw[answerred,line width=.7pt] (470.905,344.380) -- (470.953,339.047);
\draw[answerred,line width=.7pt] (470.953,339.047) -- (465.619,338.999);
\ans{469.79}{314.83}{20.00}{9}{高}
\draw[answerred,line width=.8pt] (102.000,526.667) -- (214.000,526.667) -- (152.000,467.333) -- cycle;
\draw[answerred,line width=.8pt] (321.333,526.667) -- (482.667,526.667) -- (356.667,486.667) -- cycle;
\ans{90.00}{523.33}{30.00}{9}{$\vphantom{0}A$}
\ans{216.67}{523.33}{30.00}{9}{$\vphantom{0}B$}
\ans{148.00}{454.00}{30.00}{9}{$\vphantom{0}C$}
\ans{108.67}{512.67}{30.00}{9}{$\vphantom{0}\alpha$}
\ans{194.67}{512.67}{30.00}{9}{$\vphantom{0}\beta$}
\ans{148.00}{481.33}{30.00}{9}{$\vphantom{0}\gamma$}
\ans{134.00}{528.00}{30.00}{9}{$\vphantom{0}AB$}
\ans{109.33}{490.00}{30.00}{9}{$\vphantom{0}CA$}
\ans{188.00}{490.00}{30.00}{9}{$\vphantom{0}BC$}
\ans{310.00}{523.33}{30.00}{9}{$\vphantom{0}A$}
\ans{484.67}{523.33}{30.00}{9}{$\vphantom{0}B$}
\ans{351.33}{474.00}{30.00}{9}{$\vphantom{0}C$}
\ans{338.00}{514.00}{30.00}{9}{$\vphantom{0}\alpha$}
\ans{451.33}{514.00}{30.00}{9}{$\vphantom{0}\beta$}
\ans{360.67}{496.67}{30.00}{9}{$\vphantom{0}\gamma$}
\ans{390.00}{528.00}{30.00}{9}{$\vphantom{0}AB$}
\ans{324.67}{494.00}{30.00}{9}{$\vphantom{0}CA$}
\ans{415.33}{496.67}{30.00}{9}{$\vphantom{0}BC$}
\ans{65.33}{544.67}{493.33}{8.5}{顶点：$A,B,C$；边：$AB,BC,CA$；角：$\alpha=\angle A,\ \beta=\angle B,\ \gamma=\angle C$。}
\draw[answerred,line width=.8pt] (107.333,746.667) -- (107.333,676.667) -- (177.333,746.667) -- cycle;
\draw[answerred,line width=.7pt] (107.333,740.000) -- (114.000,740.000);
\draw[answerred,line width=.7pt] (114.000,740.000) -- (114.000,746.667);
\draw[answerred,line width=.8pt] (278.000,734.667) -- (371.333,734.667) -- (324.667,711.333) -- cycle;
\draw[answerred,line width=.8pt] (433.333,746.667) -- (526.667,746.667) -- (480.000,664.667) -- cycle;
\end{scope}
\end{scope}\pgfresetboundingbox\path[use as bounding box] (62.0,585.0) rectangle (547.0,775.0);
\end{tikzpicture}}\par



\par







\QHeading{第3、4题　另一组符合条件的三角形}

\Answer{下面只给示例；大小、位置可以变化。等腰直角三角形、等腰钝角三角形和等腰锐角三角形都能画出很多个。}
\par\smallskip\noindent\begin{adjustbox}{max width=\linewidth,center}
\begin{tikzpicture}[x=.65cm,y=.65cm,font=\normalsize]
\draw[answerred,thick] (0,0)--(2,0)--(0,2)--cycle;
\draw[answerred] (0,.22)--(.22,.22)--(.22,0);
\node[below] at (1,0){等腰直角};
\begin{scope}[xshift=2.3cm]
\draw[answerred,thick] (0,0)--(3,0)--(1.5,.7)--cycle;
\node[below] at (1.5,0){等腰钝角};
\end{scope}
\begin{scope}[xshift=5cm]
\draw[answerred,thick] (0,0)--(2,0)--(1,2)--cycle;
\node[below] at (1,0){等腰锐角};
\end{scope}
\end{tikzpicture}
\end{adjustbox}\par\smallskip
\Answer{已指定底边的三角形，其对应的高必须从对面的顶点作出；同一底边的高只有一条，不能把任意垂线都当作这条高。}
